% Generated draft: requires mathematical and source review.
\documentclass[11pt]{article}
\usepackage[margin=25mm]{geometry}
\usepackage{fontspec}
\setmainfont{Latin Modern Roman}
\usepackage{amsmath,amssymb,mathtools,unicode-math}
\setmathfont{Latin Modern Math}
\usepackage{enumitem,needspace,xurl,hyperref}
\setlength{\parindent}{0pt}
\setlength{\parskip}{5pt}
\setlength{\emergencystretch}{3em}
\newcommand{\N}{\mathbb{N}}
\newcommand{\Z}{\mathbb{Z}}
\newcommand{\Q}{\mathbb{Q}}
\newcommand{\R}{\mathbb{R}}
\newcommand{\C}{\mathbb{C}}
\newcommand{\F}{\mathbb{F}}
\newcommand{\A}{\mathbb{A}}
\newcommand{\PP}{\mathbb{P}}
\newcommand{\E}{\mathbb{E}}
\newcommand{\eps}{\varepsilon}
\newcommand{\dd}{\,\mathrm{d}}
\begin{document}
\section*{UnsolvedMath ID: 801104}
\textbf{Source problem code:} catalogue entry
\subsection{Definitions and mathematical statement}
Under the standard tracial interpretation, let \(I\) be an index set, let \(\mathcal U\) be an ultrafilter on \(I\), and let \((d_i)_{i\in I}\) be positive integers satisfying \(\lim_{i\to\mathcal U}d_i=\infty\). Equip \(M_{d_i}(\C)\) with its normalized trace \(\operatorname{tr}_{d_i}\), and put
\[
\|x\|_{2,i}=\operatorname{tr}_{d_i}(x^*x)^{1/2}.
\]
The tracial matrix ultraproduct is
\[
\prod_{i\to\mathcal U}\bigl(M_{d_i}(\C),\operatorname{tr}_{d_i}\bigr)
=
\frac{\{(x_i)_{i\in I}:x_i\in M_{d_i}(\C),\ \sup_i\|x_i\|<\infty\}}
{\{(x_i)_{i\in I}:\sup_i\|x_i\|<\infty,\ \lim_{i\to\mathcal U}\|x_i\|_{2,i}=0\}}.
\]
Write \(M\) for the ambient matrix ultraproduct denoted by \(\mathbb M^{\mathcal U}\) in the source. A subfactor \(N\subseteq M\) is understood here to be a unital inclusion of von Neumann factors; finite index means that its Jones index satisfies \([M:N]<\infty\).

The supplied question asks whether every such finite-index subfactor \(N\subseteq M\) is itself isomorphic to a matrix ultraproduct. If the target is allowed to have its own index set and ultrafilter, its explicit form is
\[
[M:N]<\infty
\quad\Longrightarrow\quad
\exists\,J,\mathcal V,(e_j)_{j\in J}\quad
N\cong\prod_{j\to\mathcal V}
\bigl(M_{e_j}(\C),\operatorname{tr}_{e_j}\bigr),
\qquad \lim_{j\to\mathcal V}e_j=\infty,
\]
where \(J\) is an index set, \(\mathcal V\) is an ultrafilter on \(J\), the \(e_j\) are positive integers, and \(\cong\) denotes a normal unital \(*\)-isomorphism of von Neumann algebras.

The supplied record does not define \(\mathbb M^{\mathcal U}\), its matrix dimensions, its index set, or its ultrafilter conventions. The tracial construction and the unital Jones-index convention above are explicit standard interpretations appropriate to the supplied \(\mathrm{II}_1\)-factor context. The record also does not specify whether the target must use the same index set or the same ultrafilter. Thus the displayed implication states the unrestricted-target reading; any requirement to retain those data remains unresolved.
\subsection{Short English statement}
Must every finite-index subfactor of a tracial matrix ultraproduct itself be isomorphic to a tracial matrix ultraproduct? The supplied record leaves the ambient ultraproduct conventions and possible restrictions on the target ultrafilter unspecified.
\subsection{Sources}
\begin{itemize}
\item Local UnsolvedMath registry: permanent ID \(801104\), exact source problem code \texttt{catalogue entry}, title ``Ultraproduct \(\mathrm{II}_1\) Factors --- Problem 5.'' The supplied registry attributes the problem to Elayavalli, with proposed year 2025, and labels it \texttt{open}; this is a supplied status label, not an independently verified assessment.
\item Supplied source list: Elayavalli, \emph{50 Open Problems: Ultraproduct \(\mathrm{II}_1\) Factors} (2025), Problem 5. Supplied URL: \url{https://arxiv.org/abs/2511.20377}. The provenance records access on 2026-10-16 and extraction by \texttt{arxiv\_latex}; these are supplied provenance details, not a claim of independent source access or verification.
\end{itemize}
\end{document}
